% -*- coding: utf-8 -*-
% =============================================================================
% 小节纯切片最小示范 (例如 sec1.tex 或 sec2.tex)
% 合法环境/宏、注释规范与配置字段的权威说明见 references/tex-interface.md；
% 写作判定口径见 prompts/subagent_prompt.md。本文件注释只解释本示例的特殊点：
% 1. 本示例演示四种题型（进入/辨认/反例/证明）在同一条 probchain 中的最小写法，
%    实际切片按原文节点的实际依赖选择题环节，不得照抄本题型组合。
% 2. 本示例演示 microknowledge（实线，前置定义）与 knowledgebox（虚线，节末结论）
%    的分工：么元唯一性是学生先观察后命名的前置事实，故放实线卡。
% 3. 证明题的 proofstrategy 段只给障碍与中间目标，不泄露结论；三个小问
%    （写已知 → 结合律变形 → 收束）演示“关键论证不遗漏但支架不重复全证”。
% 4. 对话角色名示范：教师固定 $\Psi$，学生用 $\alpha$/$\beta$/$\gamma$/$\delta$ 等
%    希腊字母，严禁 Teacher/Alpha 等英文名。
% =============================================================================

\section*{I 小节标题（例如：群的概念与实例）}

% sectiondialogue：\section* 之后、首个实质题/定义之前；默认 4–8 句。
% 结构：观察 → 猜测 → 质疑/反例 → 暂时结论 → 下一任务；末句给出任务但不宣布结论。
\begin{sectiondialogue}[入口：运算的逆元是否总唯一]
\speaker{$\Psi$}{上节我们看到一些运算存在“不动的元素”。如果每个元素都有逆元，它会不会不唯一？}
\speaker{$\delta$}{一个元素当然只有一个逆元，不然怎么叫“它的”逆元？}
\speaker{$\alpha$}{“它的逆元”还没有定义。左逆元和右逆元可能不是同一个，先写出等式再谈唯一。}
\speaker{$\beta$}{我选一个具体运算表数一数，看能否造出两个不同的逆元。}
\speaker{$\Psi$}{本节任务：把“逆元唯一”从直觉变成可检验的命题，找出证明真正用到的条件。}
\end{sectiondialogue}

简要背景情境引导语（1~2 句话直接切入首题；若需说明与前题联系或研究目的，以自然语言融入题干，不要使用中括号机械标签；如背景已明确亦可直接提问）。

% 知识点多条陈列示例：必须使用 enumerate 环境，严禁散落堆砌
% \begin{enumerate}
%   \item 性质一……
%   \item 性质二……
%   \item 性质三……
% \end{enumerate}

% exampledialogue：首个入口题的引例对话（默认 3–6 句），位于 \begin{probchain}
% 之前（enumerate 内首个 \item 之前不允许出现盒子环境），完成四件事——
% 指出对象与可观察事实 → 一个角色提出可能错误的猜测 → 另一角色要求计算/分类/
% 反例 → 以尚未完成的动作结束（由题目承担该动作）。必须用注释绑定入口题 qid：
\begin{exampledialogue}[逆元个数之猜]
% dialogue: qid=concept-entry
\speaker{$\beta$}{在这个运算表里，每个元素我都找到了至少一个逆元。}
\speaker{$\gamma$}{会不会有元素同时有两个逆元？表太小，数不出来规律。}
\speaker{$\alpha$}{先把“逆元”满足的等式写出来，再用具体元素验证。}
\speaker{$\Psi$}{请大家完成下面这道题：逐个验证运算性质，再回答猜测。}
\end{exampledialogue}

% 推荐：probchain 环境自动接续上一小节的题号，无需任何 series/resume 参数
\begin{probchain}

% -----------------------------------------------------------------------------
% qid: concept-entry | 题型: 进入题 | 难度: ★☆☆
% -----------------------------------------------------------------------------
\item \qtype{进入题} 为了直观观察数集运算的结合性与特殊元素的存在现象，我们首先考察以下具体运算：
\begin{enumerate}
    \item 针对具体数集，验证运算性质：\fillin[成立]{20mm}。
    \ansspace{1.2cm}
    \item 是否满足结合律？\fillin[是]{20mm}。
    \ansspace{1.2cm}
\end{enumerate}
\begin{solution}
(a) 成立；(b) 是。
\end{solution}
\begin{teacherNote}
核心引导学生通过直观计算发现运算规律。
\end{teacherNote}

% -----------------------------------------------------------------------------
% qid: concept-identify | 题型: 辨认题 | 难度: ★☆☆
% -----------------------------------------------------------------------------
\item \qtype{辨认题} 结合上题对运算性质的检验，进一步观察其中特殊元素的存在性：\fillin[0]{15mm}。
\ansspace{1.5cm}
\begin{solution}
0
\end{solution}
\begin{teacherNote}
引导学生辨析不同系统下么元或特殊元素的表现形式。
\end{teacherNote}

% -----------------------------------------------------------------------------
% qid: concept-counterexample | 题型: 反例题 | 难度: ★★☆
% -----------------------------------------------------------------------------
\item \qtype{反例题} 为体会条件的必要性，考虑变动某一约束条件：若缺失某个条件，代数结构是否依然保持？试举出一个失效的反例并说明理由。
\ansspace{2.0cm}
\begin{solution}
例如当元素为 0 时不存在乘法逆元，导致逆元性质失效。
\end{solution}
\begin{teacherNote}
让学生通过反例深刻体会可逆性条件的必要性。
\end{teacherNote}

% -----------------------------------------------------------------------------
% [微型知识点] 一张 microknowledge 卡解决一个当前必需的定义、记号或运算。
% 保留对象、条件和必要公式，优先简短；不以固定行数删掉数学条件或缩小字体。
% 已经在前文建立的结论放入反馈或 knowledgebox，不重复包装成新定义。
% 多条性质陈列时必须使用 enumerate 环境。
% -----------------------------------------------------------------------------
\begin{microknowledge}[重要概念/性质]
\begin{enumerate}
    \item 运算结合律：对集合中任意元素 $a, b, c$，恒有 $(a \cdot b) \cdot c = a \cdot (b \cdot c)$。
    \item 么元唯一性：若存在满足 $e \cdot a = a \cdot e = a$ 的元素 $e$，则该元素必唯一。
\end{enumerate}
\end{microknowledge}

% -----------------------------------------------------------------------------
% qid: concept-proof | 题型: 证明题 | 难度: ★★★
% -----------------------------------------------------------------------------
% proofstrategy：按方法新颖程度配置（加粗引导段，不泄露待填答案，不与题干重复交代同一全证）
\textbf{证明策略：}先写出左、右逆元满足的等式，再借助结合律构造同一三元乘积的两种展开，最后比较两端收束结论。

\item \qtype{证明题} 基于上述微型知识点，我们进一步完成代数结构的严格命题推导。设运算满足结合律且存在么元，试证明逆元的唯一性：
\begin{enumerate}
    \item 明确目标：写出左逆元 $b$ 与右逆元 $c$ 各自满足的等式：\fillin[$ba=e$，$ac=e$]{25mm}。
    \ansspace{1.2cm}
    \item 凑关键步骤：考虑三元乘积表达式 $b \cdot (a \cdot c)$ 与 $(b \cdot a) \cdot c$，利用结合律展开。
    \ansspace{1.5cm}
    \item 收束结论：由此推导得出 $b = c$，即逆元唯一。
    \ansspace{1.2cm}
\end{enumerate}
\begin{solution}
由结合律 $b(ac) = (ba)c$，代入得 $b \cdot e = e \cdot c$，即 $b = c$。
\end{solution}
\begin{teacherNote}
切忌让学生一步到位证明逆元唯一，必须拆解为写出已知、结合律变形、等式收束三问。
\end{teacherNote}

\end{probchain}

% 章节末结论框（虚线）：每节必须有，仅放本节推出的结论；前置定义请用实线 microknowledge。
\begin{knowledgebox}[本节结论]
\begin{enumerate}
  \item 结论一：……
  \item 结论二：……
\end{enumerate}
\end{knowledgebox}
